\title{DeepSeekMath: Pushing the Limits of Mathematical Reasoning in Open Language Models}

\begin{abstract}
Mathematical reasoning remains a formidable task for language models, largely due to its inherently intricate and highly structured characteristics. In this work, we present DeepSeekMath~7B, which builds upon DeepSeek-Coder-Base-v1.5 7B through continued pre-training using 120B math-centric tokens collected from Common Crawl, as well as additional natural language and code content. DeepSeekMath~7B attains a notable 51.7\% on the challenging, competition-level MATH benchmark, achieving this without dependence on external toolkits or majority-vote methods. This performance rivals that of Gemini-Ultra and GPT-4. Employing a self-consistency approach with 64 generated samples, DeepSeekMath~7B reaches a score of 60.9\% on the MATH benchmark. The model’s strong proficiency in mathematical reasoning can be traced to two principal factors: firstly, the exploitation of the vast quantities of open web data via a carefully devised data selection process; secondly, the integration of Group Relative Policy Optimization (GRPO)—an adaptation of Proximal Policy Optimization (PPO)—which both augments the model’s mathematical capabilities and optimizes PPO’s memory requirements.
\end{abstract}

\section{Introduction}

Large language models (LLM) have fundamentally transformed methodologies in mathematical reasoning within the field of artificial intelligence, fostering substantial progress on tasks such as the quantitative reasoning benchmark \citep{MATH} and the geometry reasoning benchmark \citep{trinh2024solving}. Additionally, these systems have demonstrated their value as tools for supporting humans in solving sophisticated mathematical challenges \citep{tao}. Despite this, top-performing models like GPT-4 \citep{gpt4} and Gemini-Ultra \citep{gemini} remain closed-source, resulting in a notable performance gap between these and the best open-source alternatives currently available.

This work presents DeepSeekMath, a language model tailored to mathematical domains that demonstrates markedly superior mathematical proficiency compared to existing open-source counterparts, coming close to GPT-4's benchmark results on scholarly evaluations. Central to this achievement is the construction of the DeepSeekMath Corpus: an extensive, high-quality pre-training resource containing 120B mathematics-related tokens. The data is harvested from the Common Crawl (CC) via a classifier utilizing fastText \citep{joulin2016fasttext}. Initially, the classifier is trained using positive samples from OpenWebMath \citep{openwebmath} and a heterogeneous mix of unrelated web pages as negative samples. The trained classifier is then employed to identify further positive math examples within the CC, which are subsequently curated through manual annotation. This newly curated dataset is incorporated to retrain and enhance the classifier's effectiveness. Evaluation demonstrates the corpus’s substantial quality, with DeepSeekMath-Base 7B attaining 64.2\% accuracy on GSM8K \citep{gsm8k} and 36.2\% on the competitive MATH dataset \citep{MATH}, surpassing the results achieved by Minerva 540B \citep{minerva}. Moreover, given the multilingual nature of the DeepSeekMath Corpus, improved results are observed on Chinese mathematical evaluations \citep{wei2023cmath,agieval}. We posit that our methodologies in mathematical data handling serve as a valuable foundation for further academic exploration, acknowledging considerable opportunity for continued advancement.

DeepSeekMath-Base is constructed by initializing from DeepSeek-Coder-Base-v1.5 7B \citep{deepseek-coder}, motivated by findings that leveraging a code-centric pre-trained model is advantageous relative to initializing from a generic large language model. Additionally, we discover that training on mathematical tasks not only improves mathematical performance, but also yields gains on general-purpose reasoning benchmarks, such as MMLU \citep{mmlu} and BBH \citep{bbh}. This suggests that mathematical pre-training benefits broader reasoning skills beyond the mathematical domain.

Following pre-training, we further refine DeepSeekMath-Base with mathematical instruction tuning that incorporates chain-of-thought \citep{cot}, program-of-thought \citep{pot,pal}, and tool-augmented reasoning datasets \citep{tora}. The resultant DeepSeekMath-Instruct 7B model outperforms all other open 7B models and is competitive with instruction-tuned models possessing up to 70B parameters.

Moreover, we propose Group Relative Policy Optimization (GRPO), a novel reinforcement learning algorithm adapted from Proximal Policy Optimization (PPO) \citep{schulman2017proximal}. Unlike PPO, GRPO dispenses with the explicit use of a critic model, instead deriving its baseline from groupwise scores, which drastically curtails computational demands. Using just a targeted subset of English instruction tuning data, GRPO substantially improves upon the already strong DeepSeekMath-Instruct performance, both on in-domain (GSM8K: 82.9\% $\rightarrow$ 88.2\%, MATH: 46.8\% $\rightarrow$ 51.7\%) and out-of-domain tasks (for example, CMATH: 84.6\% $\rightarrow$ 88.8\%) during the RL stage. 

Additionally, we present a unified framework to contextualize several related algorithms, such as Rejection Sampling Fine-Tuning (RFT) \citep{yuan2023scaling}, Direct Preference Optimization (DPO) \citep{dpo}, PPO, and GRPO. Within this framework, we show these approaches can all be interpreted as forms of direct or simplified reinforcement learning. To dissect the critical components underlying this paradigm, we perform a broad set of experiments—including comparisons between online and offline training, outcome versus process supervision, as well as single-turn versus iterative RL strategies. Finally, we articulate the reasons behind the observed enhancement of instruction-tuned models following RL and outline prospective avenues for advancing reinforcement learning efficacy within this unified paradigm.

\subsection{Contributions}
Our work presents advancements in large-scale math pre-training, alongside investigations into reinforcement learning techniques.

\textbf{Math Pre-Training at Scale}

\begin{itemize}[topsep=0pt]
    \item We demonstrate that the openly available Common Crawl resource is a rich repository for mathematical data. Utilizing a carefully crafted data filtering pipeline, we curate the DeepSeekMath Corpus—a rigorously filtered dataset containing 120B tokens sourced from web pages focused on mathematics. This dataset is nearly seven times larger than that used in Minerva \citep{minerva} and nine times greater than OpenWebMath's recent release \citep{openwebmath}.

    \item Our DeepSeekMath-Base 7B pre-trained model attains similar results to Minerva 540B \citep{minerva}, suggesting that sheer parameter count does not solely determine mathematical reasoning proficiency. Effective training with high-quality data can enable relatively smaller models to exhibit competitive performance.

    \item Through our empirical studies on mathematical training, we observe that pre-training models on code prior to math data enhances their mathematical problem-solving skills, both when leveraging external tools and without them. This addresses the persistent query: \textit{does code training improve reasoning ability?} Our results support the notion that, at least in mathematical domains, it does.

    \item While pre-training on arXiv documents remains a widespread approach in math-related model development, our experiments reveal no significant gains across any of the mathematical benchmarks assessed in this work.
\end{itemize}

\textbf{Exploration and Analysis of Reinforcement Learning}

\begin{itemize}[topsep=0pt]
    \item We propose Group Relative Policy Optimization (GRPO), a novel reinforcement learning algorithm that eschews the use of a critic model. Instead, GRPO derives its baseline from aggregated group scores, which markedly lowers computational costs and training time in comparison to Proximal Policy Optimization (PPO).
    \item Our results indicate that GRPO substantially improves the instruction-tuning performance of our DeepSeekMath-Instruct model when relying solely on instruction-tuning datasets. In addition, we observe notable gains in the model's ability to generalize to out-of-domain tasks during reinforcement learning.
    \item We establish a unified framework for interpreting various methods—such as RFT, DPO, PPO, and GRPO—and perform an array of in-depth experiments, including analyses of online versus offline training, outcome versus process supervision, and single-turn versus iterative reinforcement learning, among others, to closely examine the critical aspects of this unified approach.
    \item Drawing from our unified perspective, we investigate the underlying mechanisms that contribute to the efficacy of reinforcement learning and outline several avenues for further enhancing reinforcement learning for LLMs.
\end{itemize}

\subsection{Summary of Evaluations and Metrics}

\begin{itemize}[topsep=0pt]
    \item \textbf{English and Chinese Mathematical Reasoning}: Our models are systematically evaluated on a variety of English and Chinese benchmarks spanning mathematical challenges from elementary school to university curricula. English datasets include GSM8K \citep{gsm8k}, MATH \citep{MATH}, SAT \citep{llemma}, OCW Courses \citep{minerva}, and MMLU-STEM \citep{mmlu}, while the Chinese benchmarks feature MGSM-zh \citep{mgsm}, CMATH \citep{wei2023cmath}, Gaokao-MathCloze \citep{agieval}, and Gaokao-MathQA \citep{agieval}. Our evaluation assesses both the models' capacities for producing independent, comprehensive solutions in text form and their proficiency in resolving problems with Python code.

    On English assessments, DeepSeekMath-Base achieves performance levels comparable to the closed-source Minerva 540B \citep{minerva}, and outperforms all existing open-source base models (such as Mistral 7B \citep{mistral} and Llemma-34B \citep{llemma}), irrespective of math pre-training, often by a notable margin. Importantly, DeepSeekMath-Base demonstrates even stronger results on Chinese evaluations, which may be attributed to our distinct approach of supplementing training data with high-quality mathematical content beyond English, in contrast to prior studies \citep{minerva,llemma}. Augmenting the model with mathematical instruction tuning and reinforcement learning, DeepSeekMath-Instruct and DeepSeekMath-RL achieve robust outcomes, surpassing 50\% accuracy on the challenging competition-level MATH benchmark for the first time within open-source initiatives.
\end{itemize}

\item \textbf{Formal Mathematics}: The DeepSeekMath-Base model is assessed on the informal-to-formal theorem proving benchmark introduced by \citep{dsp_proof}, utilizing miniF2F \citep{minif2f} as the testbed and Isabelle \citep{isabelle} as the selected proof assistant. Results indicate that DeepSeekMath-Base achieves robust few-shot performance in autoformalization.
\item \textbf{Natural Language Understanding, Reasoning, and Code}: In order to provide a holistic evaluation of the model’s capabilities in general comprehension, reasoning, and programming, DeepSeekMath-Base is benchmarked on Massive Multitask Language Understanding (MMLU) \citep{mmlu}, which consists of 57 diverse multiple-choice tasks, BIG-Bench Hard (BBH) \citep{bbh} featuring 23 multi-step reasoning tasks, as well as on code-centric datasets such as HumanEval \citep{codex} and MBPP \citep{mbpp}. Pre-training on mathematical data has been found to enhance both linguistic comprehension and reasoning abilities.

\section{Math Pre-Training}

\subsection{Data Collection and Decontamination} This section details the methodology used to create the DeepSeekMath Corpus by leveraging Common Crawl data. Figure \ref{fig:data_collect} illustrates the iterative workflow devised to amass a large-scale mathematical dataset, beginning with an initial high-quality seed collection focused on mathematics. Notably, this data gathering framework is generalizable and can be adapted for other specialized domains, such as programming.

For the seed, OpenWebMath \citep{openwebmath}, an aggregation of meticulously curated mathematical web content, is employed. Based on this seed, a fastText model \citep{joulin2016fasttext} is trained to identify and retrieve web pages that resemble those in OpenWebMath. The training dataset consists of 500,000 samples drawn at random from the seed (serving as positive instances) and 500,000 additional web pages from Common Crawl (acting as negatives). The open-source fastText implementation\footnote{\url{https://fasttext.cc}} is used, with hyperparameters set as follows: vector size of 256, a learning rate of 0.1, word n-gram length up to 3, a minimum of 3 word occurrences, and training conducted over 3 epochs. The initial Common Crawl collection is reduced in scale via deduplication based on URL similarity and near-duplicate detection, yielding a refined set of 40B HTML documents. Subsequently, mathematical pages are retrieved from this deduplicated dataset through the fastText classifier. To eliminate irrelevant or low-quality entries, pages are ranked according to the fastText model’s confidence scores, retaining only those with the highest rankings. The subset size to retain is determined by preliminary pre-training tests across the 40B, 80B, 120B, and 160B token thresholds; for the initial iteration, the top 40B tokens are selected.

Following the initial round of data collection, a significant number of mathematical web pages are left undiscovered, largely due to the insufficient diversity in the positive examples used to train the fastText model. To address this, we expand the seed corpus by incorporating additional mathematical web sources, aiming to refine the fastText model further. In this process, the entire Common Crawl is partitioned into non-overlapping domains, where each domain comprises web pages with the same base URL. We determine, for each domain, the proportion of pages that were collected in the first round. Domains in which more than 10\% of the pages were gathered are labeled as math-related (for example, \url{mathoverflow.net}).
We then proceed with manual annotation of URLs that are linked to mathematical content within these math-oriented domains (such as \url{mathoverflow.net/questions}). Any web page associated with these URLs that was not previously collected is included in the expanded seed corpus. Through this method, we acquire a broader range of positive examples and retrain the fastText model, enhancing its capability to identify mathematical data in the subsequent cycle. After four rounds of collection, we obtain 35.5M mathematical web pages, amassing a total of 120B tokens. By the fourth cycle, we observe that about 98\% of the pages had already been acquired by the third round, prompting us to halt further collection.

To prevent contamination of evaluation benchmarks, we adopt the methodology from \cite{deepseek-coder} by filtering out any web pages that feature questions or solutions from English mathematical benchmarks (such as GSM8K \citep{gsm8k} and MATH \citep{MATH}), as well as from Chinese benchmarks (like CMATH \citep{wei2023cmath} and AGIEval \citep{agieval}). The filtering process is as follows: any segment in the text containing a contiguous sequence of 10 tokens that matches any substring in these benchmarks is excluded from our mathematical training dataset. For benchmark fragments shorter than 10 tokens but of at least 3 tokens, we use exact string matching to eliminate contaminated content.

\subsection{Validation of the DeepSeekMath~Corpus Quality}

To assess the relative quality of the DeepSeekMath~Corpus, we conduct pre-training studies using the following recently published mathematical corpora as baselines:

\begin{itemize}[topsep=0pt]
    \item \textbf{MathPile} \citep{mathpile}: a compilation of materials (8.9B tokens) sourced from textbooks, Wikipedia, ProofWiki, CommonCrawl, StackExchange, and arXiv, with over 85\% originating from arXiv;
    \item \textbf{OpenWebMath} \citep{openwebmath}: a dataset extracted from CommonCrawl and filtered for mathematics content, totaling 13.6B tokens;
    \item \textbf{Proof-Pile-2} \citep{llemma}: a mathematical training set that merges OpenWebMath, AlgebraicStack (comprising 10.3B tokens of math code), and arXiv papers (28.0B tokens). For Proof-Pile-2, our experiments follow the arXiv:Web:Code sampling ratio of 2:4:1 specified by \cite{llemma}.
\end{itemize}

\subsubsection{Training Setting}
\label{sec:quality-policy}
We specialize a general-purpose pre-trained language model consisting of 1.3B parameters, architecturally aligned with DeepSeek LLMs \citep{deepseek-llm} and referred to as DeepSeek-LLM 1.3B, through additional training on mathematical corpora. Each mathematical dataset is used to independently train a model for a total of 150B tokens. All training runs utilize the lightweight and efficient HAI-LLM \citep{haillm} training platform. In accordance with DeepSeek LLMs' training methodology, we employ the AdamW optimizer \citep{adamW}, specifying $\beta_1=0.9$, $\beta_2=0.95$, and $\mathrm{weight\_decay}=0.1$. The learning rate follows a multi-phase schedule, reaching its maximum value after 2,000 warmup iterations, dropping to 31.6\% of its peak at the 80\% mark of total training steps, and further reducing to 10.0\% at the 90\% completion point. The learning rate's upper bound is fixed at 5.3e-4, with each batch comprising 4M tokens and a maximum sequence length of 4K.

\subsubsection{Evaluation Results}

\textbf{The DeepSeekMath~Corpus distinguishes itself by its high quality, broad multilingual representation, and unparalleled scale among mathematical datasets.}

\begin{itemize}[topsep=0pt]
    \item \textbf{High-quality}:
    Downstream effectiveness is assessed via eight mathematical evaluation benchmarks using few-shot chain-of-thought prompting \cite{cot}. Table \ref{tab:corpora_comparison} evidences that the model fine-tuned on the DeepSeekMath~Corpus clearly surpasses counterparts trained on other corpora. Furthermore, Figure \ref{fig:corpora_comparisons} indicates that, at the 50B token milestone (equivalent to a single epoch for Proof-Pile-2), the DeepSeekMath-trained model exceeds the performance of the one trained on Proof-Pile-2, suggesting a higher average data quality in DeepSeekMath.
    \item \textbf{Multilingual}:
    The DeepSeekMath~Corpus aggregates mathematical data in several languages, with English and Chinese being the most prevalent. Table \ref{tab:corpora_comparison} reveals that leveraging the DeepSeekMath~Corpus yields significant advancements in mathematical reasoning tasks for both English and Chinese. By contrast, many existing corpora focus predominantly on English, yielding only modest benefits or even impeding Chinese mathematical task performance.
    \item \textbf{Large-scale}:
    In terms of volume, the DeepSeekMath~Corpus significantly outweighs previously established mathematical datasets. As illustrated in Figure \ref{fig:corpora_comparisons}, training DeepSeek-LLM 1.3B on DeepSeekMath~Corpus results in a more pronounced and sustained learning improvement. The comparatively limited size of prior corpora means models trained on them experience multiple data passes and quickly saturate in performance gains.
\end{itemize}

\subsection{Training and Evaluating DeepSeekMath-Base 7B}

This section details DeepSeekMath-Base 7B, a foundational model characterized by robust mathematical reasoning capabilities. The initial weights of the model are derived from DeepSeek-Coder-Base-v1.5 7B \citep{deepseek-coder}, followed by training on a dataset comprising 500B tokens. The composition of the training data is distributed as follows: 56\% originates from the DeepSeekMath Corpus, 4\% from AlgebraicStack, 10\% from arXiv, 20\% consists of code sourced from Github, and the final 10\% includes natural language from Common Crawl in both English and Chinese. The training procedure closely mirrors the configuration described in Section \ref{sec:quality-policy}, with the exception of setting the peak learning rate to 4.2e-4 and employing a batch size of 10M tokens.

The evaluation strategy provides an in-depth analysis of DeepSeekMath-Base 7B’s proficiency in mathematical tasks. Specifically, we examine the model's ability to independently generate fully self-contained mathematical solutions, utilize tools to address mathematical problems, and formally prove theorems. Additionally, we expand our evaluation to present a broader overview of the base model’s abilities, including natural language understanding, logical reasoning, and programming competence.

\paragraph{Mathematical Problem Solving with Step-by-Step Reasoning}

DeepSeekMath-Base is assessed on mathematical problem-solving via few-shot chain-of-thought prompting \citep{cot} using eight benchmarks in both English and Chinese. These benchmarks encompass a range of quantitative reasoning tasks (such as GSM8K \citep{gsm8k}, MATH \citep{MATH}, and CMATH \citep{wei2023cmath}) as well as multiple-choice examinations (including MMLU-STEM \citep{mmlu} and Gaokao-MathQA \citep{agieval}), representing a comprehensive cross-section of mathematical domains from elementary to undergraduate complexity levels.

As detailed in Table \ref{tab:base_math_cot}, DeepSeekMath-Base 7B achieves the highest scores across all eight assessed benchmarks when compared to other open-source base models, including the widely adopted generalist model Mistral 7B \citep{mistral} and the newly introduced Llemma 34B \citep{llemma}, the latter of which benefits from dedicated math training using Proof-Pile-2 \citep{llemma}. Of particular note, DeepSeekMath-Base secures an absolute lead exceeding 10\% over all other open-source base models on the challenging MATH dataset. Furthermore, it surpasses Minerva 540B \citep{minerva}, a proprietary model that is 77 times larger, built upon PaLM \citep{palm}, and additionally refined using mathematical corpora.

\paragraph{Mathematical Problem Solving with Tool Use} We assess the models’ ability for program-assisted mathematical reasoning on the GSM8K and MATH benchmarks, utilizing a few-shot program-of-thought prompting approach \citep{pot,pal}. For each question, models are prompted to write Python code, making use of packages such as \textit{math} and \textit{sympy} to facilitate complex calculations, with their correctness evaluated based on the execution output. According to Table \ref{tab:base_math_pot_proof}, DeepSeekMath-Base 7B achieves results superior to those of the prior state-of-the-art, Llemma 34B.

\paragraph{Formal Mathematics}

Formalizing mathematical proofs can greatly improve proof correctness, reliability, and efficiency, leading to growing interest in the field. To evaluate this, we apply DeepSeekMath-Base 7B to the informal-to-formal proof generation task from \citep{dsp_proof}, which involves producing a formal proof from an informal theorem statement, its formal translation, and an informal argument. Evaluations are performed on the miniF2F benchmark \citep{minif2f}, designed for formal mathematical problems of Olympiad level, where the model is asked to construct formal proofs in Isabelle using a few-shot prompt format. Following the methodology in \cite{dsp_proof}, model outputs are initially proof sketches, and an existing automated tool, Sledgehammer \citep{sledgehammer}, is employed to complete the detailed proof steps. As indicated in Table \ref{tab:base_math_pot_proof}, DeepSeekMath-Base 7B shows robust performance in automatically formalizing mathematical proofs.

\paragraph{Natural Language Understanding, Reasoning, and Code}

We further measure performance on tasks involving language understanding via MMLU \citep{mmlu}, reasoning via BBH \citep{bbh}, and programming ability using HumanEval \citep{codex} and MBPP \citep{mbpp}. As Table \ref{tab:understanding_reasoning_code} demonstrates, DeepSeekMath-Base 7B markedly outperforms its predecessor, DeepSeek-Coder-Base-v1.5 \citep{deepseek-coder}, on MMLU and BBH, underscoring the advantages of mathematical training for language comprehension and reasoning abilities. Moreover, through the continued use of code-related tokens during training, DeepSeekMath-Base 7B successfully sustains the coding performance levels established by DeepSeek-Coder-Base-v1.5 on both programming benchmarks. Collectively, these results show that DeepSeekMath-Base 7B significantly exceeds the generalist Mistral 7B \citep{mistral} on all three tasks relating to reasoning and coding.

\section{Supervised Fine-Tuning}

\subsection{SFT Data Curation}

We assemble an instruction-tuning dataset for mathematics that spans both English and Chinese questions, incorporating a variety of mathematical domains and difficulty levels. Each problem is paired with an explanatory solution in one of several reasoning formats: chain-of-thought (CoT) \citep{cot}, program-of-thought (PoT) \citep{pot,pal}, or tool-augmented reasoning \citep{tora}. The compiled dataset comprises 776,000 training samples.

\begin{itemize}[topsep=0pt]
    \item \textbf{English mathematical datasets}: Our annotation efforts extend GSM8K and MATH tasks with tool-augmented solutions, and we integrate selected data from MathInstruct \citep{MathInstruct} and the Lila-OOD training set \citep{lila}, where the solutions employ either CoT or PoT techniques. The resulting English dataset addresses multiple mathematical areas including algebra, probability, number theory, calculus, and geometry.
    \item \textbf{Chinese mathematical datasets}: We gather Chinese K-12 mathematics problems, distributed across 76 specific subcategories (such as linear equations), with each question provided with both CoT and tool-integrated solution formats.
\end{itemize}

\subsection{Training and Evaluating DeepSeekMath-Instruct 7B}

Here, we detail the process of instruction-tuning DeepSeekMath-Instruct 7B, initialized from DeepSeekMath-Base. To maximize the model’s input capacity, training samples are concatenated randomly until a context window of 4,000 tokens is reached. The model is trained over 500 iterations, with each batch comprising 256 samples, and the learning rate is held constant at $5 \times 10^{-5}$.

We systematically assess the mathematical proficiency of the models with and without access to external tools, evaluating their performance on four quantitative reasoning benchmarks available in both English and Chinese. For comparative purposes, we position our model alongside state-of-the-art systems available at the time, including:

\begin{itemize}[topsep=0pt]
    \item \textbf{Closed-source models} considered: (1) models from the GPT series, with GPT-4 \citep{gpt4} and GPT-4 Code Interpreter~\footnote{\url{https://openai.com/blog/chatgpt-plugins\#\#code-interpreter}} representing the highest-performing variants, (2) Gemini Ultra and Pro \citep{gemini}, (3) Inflection-2 \citep{inflection-2}, (4) Grok-1~\footnote{\url{https://x.ai/model-card}}, and
    several offerings from Chinese AI companies, including (5) Baichuan-3~\footnote{\url{https://www.baichuan-ai.com}},
    (6) the newest member of the GLM-4 line~\footnote{\url{https://open.bigmodel.cn/dev/api\#glm-4}}

    from the GLM series \citep{glm}.
\end{itemize}

These models serve general applications, with most having been subjected to various alignment strategies. 
\item \textbf{Open-source models} considered here encompass both general-purpose and math-specialized architectures. The general category includes (1) DeepSeek-LLM-Chat 67B \citep{deepseek-llm}, (2) Qwen 72B \citep{qwen}, (3) SeaLLM-v2 7B \citep{seallm}, and (4) ChatGLM3 6B \citep{chatglm3}. The math-enhanced category features (5) InternLM2-Math 20B~\footnote{\url{https://github.com/InternLM/InternLM-Math}}, which is based on InternLM2, receives specialized mathematics instruction and subsequent instruction tuning; (6) Math-Shepherd-Mistral 7B, which leverages PPO training \citep{schulman2017proximal} on Mistral 7B \citep{mistral} utilizing a process-supervised reward system; (7) the WizardMath collection \citep{wizardmath}, which augments Mistral 7B and Llama-2 70B \citep{llama2} for mathematical reasoning via evolve-instruct (a variant of instruction tuning incorporating AI-generated instructions) and PPO, with tasks drawn mainly from GSM8K and MATH; (8) MetaMath 70B \citep{metamath}, a Llama-2 70B model further trained on extended GSM8K and MATH datasets; (9) ToRA 34B \cite{tora}, a variant of CodeLlama 34B tailored for mathematical reasoning with tool integration; and (10) MAmmoTH 70B \citep{MathInstruct}, which involves Llama-2 70B fine-tuned using MathInstruct through instruction tuning.
\end{itemize}

Referring to Table \ref{tab:sft_rl_math}, when evaluated without permitting the use of external tools, DeepSeekMath-Instruct 7B delivers strong stepwise reasoning abilities. On the MATH dataset, designed for competitions, it outperforms all open-source counterparts and even most closed-source systems (such as Inflection-2 and Gemini Pro) by a margin of no less than 9\% in absolute terms. This advantage holds even against models of significantly larger scale (for example, Qwen 72B) or those receiving explicit mathematical reinforcement learning (such as WizardMath-v1.1 7B). Although DeepSeekMath-Instruct's results approach those of proprietary Chinese models like GLM-4 and Baichuan-3 on MATH, it does not reach the performance of GPT-4 or Gemini Ultra.

When problem solving is evaluated with both natural language and programmatic tool use permitted, DeepSeekMath-Instruct 7B achieves close to 60\% accuracy on the MATH dataset, establishing itself as the leading open-source option. Across other benchmark tasks, the model is highly competitive with DeepSeek-LLM-Chat 67B, which represents the previous state-of-the-art but is an order of magnitude larger.

\section{Reinforcement Learning}

\subsection{Group Relative Policy Optimization} Reinforcement learning (RL) has demonstrated notable efficacy in further advancing the mathematical problem-solving proficiency of LLMs after the Supervised Fine-Tuning (SFT) phase \citep{wang2023math,wizardmath}. This section introduces our proposed RL method, Group Relative Policy Optimization (GRPO), emphasizing both its efficiency and effectiveness.

\subsubsection{From PPO to GRPO} Proximal Policy Optimization (PPO) \citep{schulman2017proximal} stands as a prominent actor-critic RL algorithm routinely employed during the RL fine-tuning phase for LLMs \citep{ouyang2022training}. PPO improves LLMs by optimizing the following surrogate loss:

\begin{equation}
\footnotesize
    \mathcal{J}_{PPO}(\theta) = \mathbb{E}{[q \sim P(Q), o \sim \pi_{\theta_{old}}(O|q)]} \frac{1}{|o|} \sum_{t=1}^{|o|} \min \left[ \frac{\pi_\theta(o_{t} | q, o_{<t})}{\pi_{\theta_{old}}(o_{t} | q, o_{<t})} A_{t}, \text{clip} \left( \frac{\pi_\theta(o_{t} | q, o_{<t})}{\pi_{\theta_{old}}(o_{t} | q, o_{<t})},

In this context, $\pi_{\theta}$ and $\pi_{\theta_{old}}$ denote the current and previous policy models, while $q$ and $o$ represent questions and outputs sampled from the question set and generated using the previous policy $\pi_{\theta_{old}}$. The hyper-parameter $\varepsilon$ is associated with the clipping mechanism in PPO to enhance training stability. The term $A_t$ stands for the advantage, which is derived using Generalized Advantage Estimation (GAE) \citep{gae}, utilizing rewards $\{r_{\ge t}\}$ in conjunction with a learned value function $V_{\psi}$. Therefore, PPO requires simultaneous training of a value function and the policy. To counteract over-optimization with respect to the reward model, a standard measure is to include a token-wise KL divergence penalty between the policy and a reference model in the reward for every token, as suggested in \citep{ouyang2022training}, specifically,

\begin{equation}
    r_{t} = r_\varphi(q, o_{\le t}) - \beta \log\frac{\pi_{\theta}(o_{t}|q, o_{<t})}{\pi_{ref}(o_{t}|q, o_{<t})},
\label{eq:PPO-reward}
\end{equation}

Here, $r_\varphi$ denotes the reward model, $\pi_{ref}$ refers to the reference model (commonly the initial SFT model), and $\beta$ serves as the KL penalty coefficient.

Given that the value function in PPO is generally another network of a scale similar to the policy model, this setup substantially increases both memory and computational demands. Moreover, during the reinforcement learning process, the value function is used as a baseline for computing the advantage to decrease variance. In large language model settings, the reward model usually assigns a reward only to the final token, potentially complicating accurate per-token value function training. To tackle these issues, as depicted in Figure \ref{fig:grpo}, we introduce Group Relative Policy Optimization (GRPO), which eliminates the need for training an auxiliary value function like PPO. Instead, GRPO employs the average reward from multiple output samples (in response to the same question) as the baseline. Specifically, for each question $q$, GRPO generates a set of outputs $\{o_1, o_2, \cdots, o_G\}$ sampled from the previous policy $\pi_{\theta_{old}}$, and optimizes the current policy to maximize the following objective:

\begin{equation}
\footnotesize
\begin{split}
    \mathcal{J}_{GRPO}(\theta) &= \mathbb{E}{[q \sim P(Q), \{o_i\}_{i=1}^G \sim \pi_{\theta_{old}}(O|q)]}  \\
    & \frac{1}{G}\sum_{i=1}^G\frac{1}{|o_i|} \sum_{t=1}^{|o_i|} \left\{ \min \left[ \frac{\pi_\theta(o_{i,t} | q, o_{i,<t})}{\pi_{\theta_{old}}(o_{i,t} | q, o_{i,<t})} \hat{A}_{i,t}, \text{clip} \left( \frac{\pi_\theta(o_{i,t} | q, o_{i,<t})}{\pi_{\theta_{old}}(o_{i,t} | q, o_{i,<t})}, 1 - \varepsilon, 1 + \varepsilon \right)  \hat{A}_{i,t} \right] - \beta \mathbb{D}_{KL}\left[\pi_{\theta} || \pi_{ref}\right]\right\} ,
\end{split}
\label{eq:GRPO-obj}
\end{equation}

Within this formulation, $\varepsilon$ and $\beta$ serve as hyper-parameters, and $\hat{A}_{i,t}$ is the group-relative advantage calculated using only the outputs within the respective group—further details are presented in subsequent subsections. By computing advantages in a group-relative fashion, GRPO naturally accommodates the comparative design of reward models, which are generally trained on paired comparisons for the same question. Furthermore, GRPO modifies the regularization strategy: instead of integrating the KL penalty directly into the reward, it incorporates the KL divergence between the policy and the reference policy as a loss regularization term, thus streamlining the estimation of $\hat{A}_{i,t}$. Notably, in contrast to the KL term in (\ref{eq:PPO-reward}), we employ the following unbiased estimator for the KL divergence \citep

\begin{algorithm}[t]
  \small
  \caption{Iterative Group Relative Policy Optimization}
  \textbf{Input} initial policy model $\pi_{\theta_{\text{init}}}$; reward models $r_\varphi$; task prompts $\mathcal{D}$; 
  hyperparameters $\varepsilon$, $\beta$, $\mu$
  \begin{algorithmic}[1]
    \State Initialize policy model: $\pi_\theta \leftarrow \pi_{\theta_{\text{init}}}$
    \For{iteration = 1, \dots, I}
       \State Set reference model $\pi_{ref} \leftarrow \pi_{\theta}$
      \For{step = 1, \dots, M}
      \State Randomly sample mini-batch $\mathcal{D}_b$ from $\mathcal{D}$
      \State Assign previous policy model: $\pi_{\theta_{old}} \leftarrow \pi_{\theta}$ 
      \State For each question $q \in \mathcal{D}_b$, generate $G$ candidate responses $\{o_i\}_{i=1}^G \sim \pi_{\theta_{old}} (\cdot \mid q) $
      \State Evaluate sampled responses with reward model $r_\varphi$ to obtain scores $\{r_i\}_{i=1}^{G}$
      \State Estimate group-relative advantages $\hat{A}_{i,t}$ for every $t$-th token in $o_i$.
      \For{GRPO iteration = 1, \dots, $\mu$}
        \State Update $\pi_\theta$ by optimizing the GRPO objective in Equation \ref{eq:GC-GRPO}
      \EndFor
    \EndFor
    \State Continue training $r_\varphi$ using a replay buffer mechanism.
    \EndFor
  \end{algorithmic}
  \textbf{Output} $\pi_\theta$
  \label{alg:iter-grpo}
\end{algorithm}

\subsubsection{Outcome Supervision RL with GRPO} To specify, for any given question $q$, a set of $G$ outputs $\{o_1, o_2, \cdots, o_G\}$ are produced by sampling from the current policy model $\pi_{\theta_{old}}$. These outputs are evaluated by a reward model, resulting in a corresponding reward vector $\mathbf{r} = \{r_1, r_2, \cdots, r_G\}$. Each reward is standardized by subtracting the group mean and dividing by the group standard deviation. Under outcome supervision, the standardized reward for output $o_i$ is assigned to all its tokens as the advantage $\hat{A}_{i, t}$, namely, $\hat{A}_{i, t} = \widetilde{r}_i = \frac{r_i- {\rm mean}(\mathbf{r})}{{\rm std}(\mathbf{r})}$, and the policy model is updated by maximizing the objective in equation (\ref{eq:GRPO-obj}).

\subsubsection{Process Supervision RL with GRPO} Since outcome supervision restricts feedback to only the terminal output, it may not provide sufficient or granular guidance for challenging reasoning problems. Inspired by \cite{wang2023math}, we also implement process supervision, which supplies rewards at every intermediate reasoning step. Concretely, with question $q$ and $G$ sampled outputs $\{o_1, o_2, \cdots, o_G\}$, a stepwise reward model rates each reasoning segment, producing: $\mathbf{R} = \{ \{r_1^{index(1)},\cdots,r_1^{index(K_1)}\}, \cdots,  \{r_G^{index(1)},\cdots,r_G^{index(K_G)}\} \}$, where $index(j)$ denotes the token at the conclusion of the $j$-th step and $K_i$ is the number of reasoning steps for output $o_i$. The stepwise rewards are normalized in the same fashion: $\widetilde{r}_i^{index(j)} = \frac{r_i^{index(j)} - {\rm mean}(\mathbf{R})}{{\rm std}(\mathbf{R})}$. For process supervision, the advantage at each token $t$ is computed as the sum of normalized rewards for steps at or after $t$, i.e., $\hat{A}_{i, t} = \sum_{index(j) \ge t} \widetilde{r}_i^{index(j)}$, and the model is optimized according to equation (\ref{eq:GRPO-obj

\subsubsection{Iterative RL with GRPO} As reinforcement learning training advances, reliance on a prior reward model often proves inadequate for guiding the continually evolving policy model. To address this limitation, we investigate iterative RL using GRPO. According to Algorithm \ref{alg:iter-grpo}, the iterative GRPO approach entails generating fresh training data for the reward model via new samples collected from the current policy model. We incrementally update the previous reward model using a replay strategy, where 10\% of samples are drawn from past data. Next, we set the policy model as the new reference model and proceed to further train it in conjunction with the updated reward model.

\subsection{Training and Evaluating DeepSeekMath-RL}

Our reinforcement learning procedure employs DeepSeekMath-Instruct 7B as its foundation. The RL training corpus is drawn from chain-of-thought-style GSM8K and MATH questions within the SFT dataset, amounting to approximately 144,000 items. Other SFT data are omitted to better isolate RL's effects on benchmarks that do not contribute training data during this phase. The assembly of training data for the reward models adheres to the methodology from \citep{wang2023math}. We initialize the reward model using DeepSeekMath-Base 7B with a 2e-5 learning rate. For GRPO optimization, the policy model employs a learning rate of 1e-6 and a KL coefficient set to 0.04. Each prompt yields $64$ sampled completions. Sequence length and training batch size are each fixed at 1024. During each iteration, the policy model receives a single parameter update post-exploration. DeepSeekMath-RL 7B is assessed with the same evaluation benchmarks used for DeepSeekMath-Instruct 7B. Here, GSM8K and MATH chain-of-thought problems are classified as in-domain, while all remaining evaluation datasets are considered out-of-domain.

Table \ref{tab:sft_rl_math} presents comparative results for both open-source and proprietary models utilizing chain-of-thought and tool-augmented reasoning strategies across English and Chinese evaluation tasks. Our key observations are: (1) DeepSeekMath-RL 7B achieves accuracy rates of 88.2\% on GSM8K and 51.7\% on MATH with chain-of-thought reasoning, outperforming all open-source models from 7B to 70B parameter sizes, and exceeding the performance of most closed-source systems. (2) Notably, DeepSeekMath-RL 7B is exclusively fine-tuned on chain-of-thought-style instructions from GSM8K and MATH, with DeepSeekMath-Instruct 7B serving as its initialization. Despite this limited instruction set, DeepSeekMath-RL 7B demonstrates substantial improvements over DeepSeekMath-Instruct 7B in all evaluation criteria, underscoring the advantages conferred by reinforcement learning.

\section{Discussion}
In the following section, we discuss the insights gained from our pre-training and reinforcement learning experiments.

\subsection{Lessons Learnt in Pre-Training}

We begin by reporting on our experiences from the pre-training phase. Unless stated otherwise, the experimental conditions follow those described in Section \ref{sec:quality-policy}. Importantly, any mention of the DeepSeekMath Corpus within this section pertains to a dataset comprising 89B tokens, obtained during the second iteration of data collection.

\subsubsection{Code Training Benefits Mathematical Reasoning}

There exists a widespread, though not yet conclusively proven, assumption that training on code data enhances a model’s reasoning capabilities. We contribute to this discussion by providing empirical evidence specifically within the context of mathematics, showing that code-oriented pre-training boosts a model's capacity for mathematical reasoning, regardless of whether external tools are involved.

To evaluate the effect of code-centric training on mathematical reasoning, we compared models trained under two distinct regimens: two-stage training and single-stage training.

\textbf{Two-Stage Training}

\begin{itemize}[topsep=0pt]
    \item \textbf{Code Training for 400B Tokens $\rightarrow$ Math Training for 150B Tokens}: DeepSeek-LLM 1.3B undergoes pre-training with 400B code tokens, after which it is further trained on an additional 150B math tokens.
    \item \textbf{General Training for 400B Tokens $\rightarrow$ Math Training for 150B Tokens}: For comparison, we replace the code tokens in the initial training phase with general tokens sampled from a comprehensive corpus curated by DeepSeek-AI. This experiment is designed to assess whether code-specific pre-training offers unique benefits over general token pre-training for the subsequent development of mathematical reasoning skills.
\end{itemize}

\textbf{One-Stage Training}

\begin{itemize}[topsep=0pt]
    \item \textbf{Math Training for 150B Tokens}: DeepSeek-LLM 1.3B is exclusively trained on 150B math tokens.
    \item \textbf{Training on a mixture of 400B Code Tokens and 150B Math Tokens}: Since introducing math training after code pre-training tends to reduce the model's coding capability, we examine whether training on a combined set of code and math tokens in a single phase might preserve or improve both coding and mathematical reasoning, and potentially prevent the problem of catastrophic forgetting.
\end{itemize}

\paragraph{Results}
The outcomes across these training regimens are presented in Table \ref{tab:code-math} and Table \ref{tab:code-math-general-eval}.

Across both the sequential (two-stage) and concurrent (one-stage) approaches, code pre-training consistently enhances program-assisted mathematical reasoning. Specifically, Table \ref{tab:code-math} shows that in the two-stage paradigm, code token pre-training alone delivers substantial gains in Python-based solution accuracy on GSM8K and MATH datasets; these gains are further improved when followed by dedicated math training. Notably, in the one-stage setting, jointly training on both code and math tokens prevents the catastrophic forgetting observed with the sequential method and simultaneously supports both coding ability (Table \ref{tab:code-math-general-eval}) and program-aided math reasoning (Table \ref{tab:code-math}).

Furthermore, code-based training aids mathematical reasoning in scenarios where tools are not employed. In the two-stage setup, beginning with code tokens leads to appreciable improvements even before specialized math training, and it accelerates learning during subsequent math training, ultimately yielding the best performance. However, when code and math tokens are combined during a single training phase, the model’s mathematical reasoning performance without tool assistance diminishes. This phenomenon may be attributed to the limited capacity of DeepSeek-LLM 1.3B, which appears insufficient to thoroughly learn both coding and mathematical competencies simultaneously.

\subsubsection{ArXiv Papers Appear Unhelpful for Enhancing Mathematical Reasoning}

While arXiv papers are a prevalent resource in pre-training datasets for mathematics \citep{minerva,gpt-f,llemma,mathpile}, there is limited empirical assessment regarding their effectiveness in boosting mathematical reasoning skills. Somewhat unexpectedly, our experimental evidence suggests that arXiv papers do not facilitate improvements in mathematical reasoning. We conduct investigations using models of varying sizes, specifically DeepSeek-LLM 1.3B and DeepSeek-Coder-Base-v1.5 7B \citep{deepseek-coder}, pre-trained on different versions of the arXiv corpus processed through distinct cleaning pipelines:

\begin{itemize}[topsep=0pt]
    \item \textbf{MathPile} \citep{mathpile}: a dataset comprising 8.9B tokens, curated via targeted cleaning and filtering strategies, with over 85\% consisting of scientific documents from arXiv;
    \item \textbf{ArXiv-RedPajama} \citep{redpajama}: a comprehensive collection of arXiv LaTeX files, from which preambles, comments, macros, and references have been stripped, yielding 28.0B tokens.
\end{itemize}

For these experiments, DeepSeek-LLM 1.3B was trained on 150B tokens, while DeepSeek-Coder-Base-v1.5 7B was trained on 40B tokens from each arXiv source. The outcomes indicate that exposure to arXiv corpora does not advance mathematical reasoning; when trained exclusively on arXiv data, both models show no observable gains, and in certain cases, their performance declines on a range of mathematical benchmarks with varying levels of challenge as evaluated in our work. These assessments encompass quantitative reasoning tasks such as GSM8K and MATH (see Table \ref{tab:arxiv-cot}), multiple-choice tests exemplified by MMLU-STEM (see Table \ref{tab:arxiv-cot}), and tasks focused on formal mathematics such as miniF2F (see Table \ref{tab:arxiv_atp}).

Nonetheless, it is important to recognize the provisional nature of these findings. We have not yet explored:

\begin{itemize}[topsep=0pt]
    \item How arXiv data influences more specialized mathematical tasks omitted from our current analyses, for instance, the conversion of formal mathematical language into more informal explanations (informalization);
    \item The impact of integrating arXiv data with other data modalities;
    \item Potential advantages that might arise from leveraging arXiv data in models of even larger scale.
\end{itemize}

Given these limitations, additional research is necessary, which we defer to subsequent studies.

\subsection{Insights from Reinforcement Learning}
\subsubsection{Towards a Unified Framework} This section introduces a generalized framework for examining diverse training strategies, including SFT, RFT, DPO, PPO, and GRPO, alongside experimental analyses of various components within this unified perspective. The general expression for the parameter gradient $\theta$ associated with a given training algorithm can be represented as:

\begin{equation}
    \nabla_{\theta}\mathcal{J}_{\textcolor{red}{\mathcal{A}}}(\theta) = \mathbb{E}[\underbrace{(q,o) \sim \textcolor{red}{\mathcal{D}}}_{Data \ Source}]\left( \frac{1}{|o|} \sum_{t=1}^{|o|}  \underbrace{GC_{{\mathcal{A}}}(q, o, t, \textcolor{red}{\pi_{{rf}}})}_{Gradient \ Coefficient}  \nabla_{\theta}\log \pi_{\theta}(o_t | q, o_{<t})\right).
\label{eq:objective}
\end{equation}

This framework consists of three fundamental elements: 1) the \textit{Data Source $\mathcal{D}$}, which supplies the inputs for training; 2) the \textit{Reward Function $\pi_{{rf}}$}, providing the evaluative signal that guides learning; and 3) the \textit{Algorithm $\mathcal{A}$}, which interprets both the training examples and reward feedback to yield the gradient coefficient $GC$, governing the strength of either penalization or reinforcement for each example. Utilizing this general structure, we review several major algorithms:

\begin{itemize}
    \item  \textbf{Supervised Fine-tuning (SFT)}: This approach adapts a pretrained model by retraining it on curated human-annotated SFT datasets.
    \item \textbf{Rejection Sampling Fine-tuning (RFT)}: Building upon SFT, RFT additionally fine-tunes the model using a subset of outputs generated by the SFT model for the same queries, but only includes those outputs that pass a correctness-based filtering process.
    \item \textbf{Direct Preference Optimization (DPO)}: DPO continues the optimization of the SFT model through training on specially augmented samples, applying a pairwise loss function that reflects preferences between outputs.
    \item \textbf{Online Rejection Sampling Fine-tuning (Online RFT)}: Distinct from traditional RFT, Online RFT starts from an SFT-initialized policy and further trains it with outputs dynamically sampled from the current version of the policy, ensuring continual adaptation.
    \item \textbf{PPO/GRPO}: These methods also initialize policies with the SFT model, but rely on reinforcement learning by collecting outputs from the live policy and applying policy optimization techniques for further improvement.
\end{itemize}

The components of these methods are outlined in Table \ref{tab:data_policy}. For an in-depth explanation of the derivation procedures, see Appendix \ref{app:analysis-rl}.

\paragraph{Observation about Data Source} We categorize data sources into two main types: online sampling and offline sampling. Online sampling involves collecting training data directly from the real-time exploration of the current policy model, whereas offline sampling makes use of data sampled from the initial SFT model. Among the methods discussed, RFT and DPO adopt the offline approach, while Online RFT and GRPO rely on the online approach.

Figure \ref{fig:rl-analysis} illustrates that Online RFT consistently surpasses RFT on both evaluated benchmarks. While the two methods show similar performance during the initial training phase, Online RFT ultimately achieves a notable lead in later stages, underscoring the efficacy of online training. This trend is expected: early on, the actor closely resembles the SFT model, and the sampled trajectories show minimal deviation. Over time, however, the actor's behavior diverges more substantially, and leveraging freshly generated data confers increasingly significant benefits.

\paragraph{Observation about Gradient Coefficient} In these algorithms, the reward signal is mapped to a gradient coefficient that guides model parameter updates. We distinguish two types of reward functions in our experiments: `Rule' and `Model'. The `Rule' function assesses response quality by evaluating the correctness of the answer, whereas the `Model' function relies on a separately trained reward model that scores responses, with its training data itself being derived from rule-based judgments. A notable distinction between GRPO and Online RFT, as revealed by Equations \ref{eq:GC-RFT} and \ref{eq:GC-GRPO}, is GRPO's dynamic adjustment of the gradient coefficient according to the reward assigned by the reward model. This facilitates the provision of varied reinforcement or penalization, proportionate to response quality. Conversely, Online RFT does not employ this adaptive mechanism: it omits penalties for incorrect answers and applies the same degree of reinforcement to all responses deemed correct.

As illustrated in Figure \ref{fig:rl-analysis}, GRPO demonstrates superior performance over online RFT, underscoring the effectiveness of adjusting the coefficients for positive and negative gradients. Moreover, GRPO+PS outperforms GRPO+OS, which underscores the advantage of employing more granular, step-aware gradient adjustments. Additionally, we examine the impact of iterative reinforcement learning by performing two consecutive iterations in our experimental setup. Figure \ref{fig:iter-rl} reveals that employing iterative RL leads to notable performance improvements, particularly noticeable after the first iteration.

\subsubsection{Why RL Works?} In this study, reinforcement learning is carried out on a selected subset of the instruction tuning dataset, and it results in marked improvements when applied to the instruction-tuned model. To shed further light on the efficacy of reinforcement learning, we compare the Pass@K and Maj@K accuracies between the Instruct and RL models across two evaluation benchmarks. As illustrated in Figure \ref{fig:maj-pass}, reinforcement learning yields gains in Maj@K but leaves Pass@K largely unchanged. This suggests that RL enhances aggregate model performance by producing a more resilient output distribution; specifically, \textbf{the improvements are likely due to an increased likelihood of obtaining the correct answer among the TopK outputs, rather than due to strengthening basic model competencies.} Along similar lines, \citep{wang2023making} describe a \textbf{misalignment problem} for reasoning tasks in SFT models, and demonstrate that SFT model reasoning can be improved using various preference alignment techniques \citep{yuan2023rrhf,song2023preference,wang2023making}.

\subsubsection{How to Achieve More Effective RL?}
\label{sec:effec_rl}
Our results show that reinforcement learning delivers strong gains in tasks focused on mathematical reasoning. Furthermore, we propose a unified framework that contextualizes different prominent training strategies under the broader umbrella of direct or simplified RL methods. According to Equation \ref{eq:objective}, three primary components structure this paradigm: Data Source, Algorithm, and Reward Function. We outline several promising directions for these three aspects.

\paragraph{Data Source} The data source forms the foundation for all training approaches. In RL contexts, the data source is typically defined as the unlabeled question prompts, accompanied by responses generated through sampling from the policy model. For this work, we restricted ourselves to questions from the instruction tuning phase and leveraged basic nucleus sampling to generate outputs. We hypothesize that this constraint may explain why our RL approach enhances only Maj@K performance. In future work, we plan to expand our RL pipeline to include question prompts from broader, out-of-distribution domains and employ \textbf{advanced sampling (decoding) strategies}, such as those using tree-search algorithms \citep{yao2023tree}. Furthermore, integrating \textbf{efficient inference techniques} \citep{xia-etal-2023-speculative,leviathan2023fast,kwon2023efficient,xia2024unlocking}—which can substantially improve exploration efficiency for policy models—will also be pivotal.

\paragraph{Algorithms} Algorithms utilize both the available data and the received reward signal to compute the gradient coefficient for updating the model’s parameters. According to Equation \ref{eq:objective}, prevailing approaches essentially place full \textbf{TRUST} in the feedback provided by the reward function, using it to adjust the conditional probability assigned to specific tokens. Nonetheless, such trust cannot guarantee the continual reliability of the reward signal, particularly in settings characterized by high complexity. A case in point is the PRM800K datasets \citep{lightman2023let}, which, despite rigorous annotation by trained experts, still exhibit an error rate of about 20\% in their labels\footnote{\url{https://github.com/openai/prm800k/issues/12\#issuecomment-1728491852}}. To address these challenges, we aim to investigate reinforcement learning algorithms that maintain effectiveness even when the reward signals are noisy. We posit that \textbf{WEAK-TO-STRONG} \citep{burns2023weak} alignment strategies have the potential to substantially reshape learning algorithm design.

\paragraph{Reward Function} The reward function serves as the principal provider of training signals during learning. In reinforcement learning, this typically takes the form of a neural reward model. We identify three critical avenues for improving reward models: 1) \textbf{Improving the generalization capacity of the reward model.} To be truly effective, reward models must be capable of handling out-of-distribution problems and more sophisticated decoding outcomes; failure to generalize may result in reinforcement learning merely preserving LLMs’ output distributions without yielding genuine improvements; 2) \textbf{Representing reward model uncertainty.} Capturing and leveraging uncertainty could function as a crucial connector between imperfect reward models and weak-to-strong learning frameworks; 3) \textbf{Constructing high-fidelity process reward models efficiently} so that fine-grained training signals can guide complex reasoning tasks \citep{lightman2023let,wang2023math}.

\section{Conclusion, Limitation, and Future Work} This paper introduces DeepSeekMath, an open-source model that surpasses all other open-source counterparts on the competition-grade MATH benchmark, and achieves performance close to that of proprietary systems. Built upon DeepSeek-Coder-v1.5 7B, DeepSeekMath undergoes a prolonged continual training regimen over 500B tokens, notably incorporating 120B mathematics-focused tokens from Common Crawl in the training set. Our comprehensive ablation analysis demonstrates that web-derived content provides valuable, high-quality mathematical data, whereas the contribution of arXiv-sourced material appears less substantial than anticipated. Furthermore, we propose Group Relative Policy Optimization (GRPO)—an extension of Proximal Policy Optimization (PPO)—which facilitates enhancements in mathematical reasoning with reduced memory overhead. Experimental findings indicate GRPO’s effectiveness remains even as DeepSeekMath-Instruct 7B attains robust benchmark scores. We also propose a unified framework for interpreting a family of reinforcement learning methods and outline promising avenues to improve reinforcement learning efficiency.

Despite DeepSeekMath's notable results on quantitative reasoning tasks, its performance on geometry and theorem-proving tasks lags behind proprietary models. For example, during preliminary evaluations, the model was unable to successfully address problems related to triangles and ellipses, pointing to potential biases in both the pre-training and fine-tuning data selection. Furthermore, due to its smaller model size, DeepSeekMath underperforms compared to GPT-4 in few-shot scenarios—GPT-4 benefits from few-shot prompts, whereas DeepSeekMath's results do not show improvement over zero-shot trials. Moving forward, we plan to enhance our data selection and engineering processes to curate higher-quality pre-training corpora. We also intend to pursue the prospective directions outlined in Section \ref{sec:effec_rl} for advancing reinforcement learning approaches in LLM development.

\end{document}